\subsection{Critical coordinate control and the cap}\label{sec:critical-annulus}

\begin{proposition}[A short critical shell cut]\label{prop:neck-critical-shell}\label{prop:critical-annulus}
Suppose a chart on \(A(1/8,8)\) in the physical source already
has the complete marking of \eqref{eq:anchored-marks}, with both
marks off the closed trim \(K_i\). Assume its constant
rescaling \(m_i\) satisfies
\begin{equation}\label{eq:neck-critical-shell}
 \|m_i-\delta_{\rm Euc}\|_\infty\le\epsilon_*,\qquad
 \int_{A(1/8,8)}|\Rm_{m_i}|^2\,dV_{m_i}\to0,\qquad
 \|D T_i\|_{2,dz}\to0,\quad T_i^a=\Delta_{m_i}z^a .
\end{equation}
Here \(\epsilon_*>0\) is fixed sufficiently small on this annulus
before taking the source limit. In exact harmonic coordinates the
last condition is automatic. Such data are impossible on a fixed
nonstandard \(M\) at its positive-class infimum. Thus the wide
data in Proposition~\ref{prop:anchored-marking} together with
\eqref{eq:neck-critical-shell} on the central chart also give a
contradiction. The same implication holds with uniform and strong
\(W^{1,4}\) convergence to Euclidean in the supplied central chart
in place of the fixed small distortion and tension conditions.
\end{proposition}

\begin{proof}
Write \(v=m-\delta_{\rm Euc}\), \(\eta=\|v\|_\infty\) and
\(T^a=\Delta_m z^a\). The reduced metric equation has the form
\[
 m^{ab}\partial_{ab}v_{jk}
 =-2\Ric_{jk}-m_{ja}\partial_kT^a-m_{ka}\partial_jT^a
   +O(|Dv|^2),
\]
using \(T^a=-m^{bc}\Gamma^a_{bc}\)
\href{https://www.numdam.org/item/10.24033/asens.1405.pdf}{(DeTurck--Kazdan,
Lemma~4.1)}. For a compactly supported cutoff \(\chi\) set
\(X=\|\chi^2D^2v\|_2\), \(G=\|\chi Dv\|_4\). Integration by
parts against \(v\) in \(\int \chi^4|Dv|^4\) gives
\(G^2\le C\eta X+C_\chi\eta^2\). Apply the Euclidean
\(\|D^2f\|_2=\|\Delta f\|_2\) identity to \(f=\chi^2v\).
Freezing the principal coefficient costs \(C\eta X\), the
quadratic term costs \(CG^2\), and the commutators cost at most
\(C_\chi(G+\eta)\). Substituting the inequality for \(G^2\)
and absorbing with \(\eta\le\epsilon_*\) yields
\[
 X\le C_\chi(\|\chi^2\Ric\|_2+\|\chi^2 DT\|_2+\eta),\qquad
 G^2\le C_\chi\eta(\|\chi^2\Ric\|_2+\|\chi^2 DT\|_2+\eta).
\]
These Euclidean component norms are comparable to metric norms
for the indicated tensor sources. Use nested cutoffs protecting
the closed trim. Rellich compactness and these estimates give a
weak \(W^{2,2}\) limit \(h\), with strong \(W^{1,s}\) convergence
for \(s<4\) and almost-everywhere convergence. The quadratic
first-derivative terms converge in distributions by strong
\(W^{1,2}\). The bounded principal coefficients pass to the
limit against a fixed \(L^2\) test by dominated convergence.
The limit has zero right-hand curvature and differentiated
tension in the reduced equation, and zero full curvature in
distributions by the all-lowered curvature formula.

For interior regularity write the reduced equation on a ball as
\(h^{ab}\partial_{ab}h+C(h,Dh)Dh=0\), with drift coefficients
bounded by \(C_0|Dh|\). This is a small perturbation of the
Dirichlet Laplacian on \(W^{2,s}\cap W^{1,s}_0\), for \(s=2,3\).
The drift is small in \(L^4\) by the estimate above and is bounded
on that operator domain by the embedding into gradients in
\(L^{4s/(4-s)}\). A cutoff times \(h\) gives a forcing in
\(L^3\). Solving there and using uniqueness at exponent two
gives \(h\in W^{2,3}_{\rm loc}\). Its gradient lies in \(L^{12}\)
and the principal coefficients are continuous. Freezing them on
smaller balls gives \(W^{2,6}_{\rm loc}\), followed by ordinary
interior bootstrapping to a smooth flat metric. In the
approximately harmonic case a limiting constant coordinate
tension is compatible with this argument, which uses its derivative.
Subtract the reduced equations for \(m_i\) and \(h\), and localize
again. The principal perturbation and the drift of size
\(C(|Dm_i|+|Dh|)\) are absorbed at exponent two. The other terms
tend to zero in \(L^2\), including coefficient differences against
the smooth jets of \(h\), the curvature and tension sources, and
the cutoff commutators. For a drift commutator use the \(L^4\)
bound on the drift and strong \(L^4\) convergence of \(m_i-h\).
This gives strong \(W^{2,2}\) and hence strong \(W^{1,4}\)
convergence on smaller interiors.

The flat metric on the simply connected annulus \(A(1/4,4)\)
has an oriented developing map \(F\), obtained by integrating a
parallel orthonormal coframe. Since \((DF)^tDF=h\), the gradient
is uniformly within \(C\epsilon_*\) of \(\SO(4)\). The geometric
rigidity estimate for the fixed Lipschitz annulus gives a
rotation \(Q\) with \(\|DF-Q\|_2\le C\epsilon_*\)
\cite[Theorem 3.1]{FrieseckeJamesMueller2002Rigidity}.
After translation, Poincar\'e and the uniform derivative bound
make \(F-Qz\) uniformly small on the closed \(A(1/3,3)\) as
\(\epsilon_*\downarrow0\). In fact a fixed nonzero supremum
there would give a fixed lower \(L^2\) bound on an interior ball.
The straight homotopy gives degree one on \(A(1/3,3)\) for
each target point in \(A(1/2,2)\). The positive Jacobian makes
its preimage unique there. The pulled-back unit sphere separates
the original trim faces, since every path across the trim must
traverse the two intermediate radii \(1/3\) and \(3\), at which
the developed radii are below and above one, respectively. In
particular this change of coordinates retains both complete marks.

On the developed \(A(1/2,2)\) the metrics tend strongly to
Euclidean in \(W^{1,4}\) and are uniformly elliptic. The same
properties already hold on a central annulus in the alternative
first-order hypothesis. Apply the curvature comparison of
Theorem~\ref{thm:caps} with a fixed short splice cutoff, leaving
an exactly flat strip about the new cut. Uniform ellipticity and
the product rule for the all-lowered curvature tensor bound the
\(L^2\) error in the chiral densities by a constant times
\[
 \|\Rm_{m_i}\|_2+\|Dm_i\|_4^2+\|Dm_i\|_2+
 \|m_i-\delta_{\rm Euc}\|_2=o(1)
\]
in the developed chart. The projection change is bounded by the
first term, so uniform convergence after development is unnecessary.
The estimate holds along the interpolation on the closed source
and on each splice separately. A conformal multiplier on the
flat strip gives the exact round product at zero chiral cost.

The splice path remains positive. If a continuous smooth metric
path on the closed type beginning at a positive class first lost
that sign, the zero first eigenvalue of the conformal Laplacian
would give a scalar-flat representative by a positive
eigenfunction. Equation~\eqref{eq:signature} and Gauss--Bonnet
then require \(E_+\ge8\pi^2\), whereas the whole path here has
energy \(E_i+o(1)<8\pi^2\). Stretch the exact product to any
finite length on the same type by increasing its radial arclength
coefficient with a bump supported in the product. That path retains
the same chiral energy, so the Yamabe constants of the stretched
metrics are bounded below by a fixed positive constant, by
\eqref{eq:budgets} in their unit-volume Yamabe representatives.
For each fixed source index any compactly supported test on one
side with an attached half-cylinder extends by zero to a long enough
stretched source and inherits that inequality. Compactify the end
conformally using \(r=e^{-t}\), so \(r^2(dt^2+\gamma)\) is smooth
and flat at the added point. The vanishing point capacity in
dimension four transports the same Yamabe bound to the capped
side by approximation with compactly supported punctured tests.
The resulting caps are spherical and conformally flat. The
per-side nonnegative upper allowances \(\gamma_{i,\pm}\) above the
actual raw energies are \(o(1)\). Both raw masses are at least
\(\theta E_i\), and for all large \(i\)
\[
 E_i+\gamma_{i,+}+\gamma_{i,-}<8\pi^2,\qquad
 E_i-I^+(M)+\max\{\gamma_{i,+},\gamma_{i,-}\}<4\pi^2\theta .
\]
Theorem~\ref{thm:nonsplitting} contradicts these two raw masses.
All normalized windows and splice profiles are fixed in advance.
The auxiliary stretch for testing positivity is performed on an
exact product after the splice at each fixed index.
\end{proof}

\begin{corollary}[Anchored comparison at separated physical scales]
\label{cor:anchored-annulus}
Under the hypotheses of Proposition~\ref{prop:anchored-marking}, assume
also the small-distortion and harmonic or differentiated tension bounds
of Proposition~\ref{prop:neck-critical-shell} on the central annulus.
Such a sequence cannot occur on a nonstandard type minimizing in the
positive subcritical band.
\end{corollary}
\begin{proof}
The first proposition places the whole inner charged ball and the
whole median exterior on opposite sides of the buffered trim, at the
same physical scale. The second constructs the marked cut there.
The discarded side has a fixed positive charge while the minimizing
excess and the two splice costs tend to zero.
\end{proof}
